\section{高斯差分隐私机制与参数精确校准}
\label{sec:gaussian}

\subsection{高斯机制的数学定义与特性}
高斯机制（Gaussian Mechanism）是差分隐私理论中处理实数向量查询的核心原语~\cite{dwork2006calibrating, dwork2014algorithmic}。设 $f: \mathcal{D} \to \mathbb{R}^d$ 为一个确定性向量查询函数，其全局 $L_2$ 敏感度定义为：
\begin{equation}
  \Delta_2(f) = \max_{D \sim D'} \|f(D) - f(D')\|_2
\end{equation}
高斯机制通过向真实向量 $f(D)$ 叠加各向同性多元高斯随机变量实现扰动：
\begin{equation}
  \mathcal{M}_{\text{Gauss}}(D) = f(D) + \mathcal{N}\left(0, \sigma^2 \mathbf{I}_d\right)
\end{equation}
高斯机制相较于拉普拉斯机制的关键理论优势在于：其每坐标噪声尺度由 $L_2$ 敏感度控制；总噪声范数通常仍随 $\sqrt d$ 增长，且在多重组合与长周期发布场景下，具有优良的 Rényi 散度可加性。

\subsection{经典高斯校准及其理论局限}
在早期差分隐私文献中，为了保证算法满足 $(\varepsilon, \delta)$-DP，通常采用基于切诺夫界的经典闭式近似公式：
\begin{equation}
  \label{eq:classical_gaussian}
  \sigma_{\text{classical}} = \frac{\Delta_2(f) \sqrt{2 \ln\left(\frac{1.25}{\delta}\right)}}{\varepsilon}
\end{equation}
然而，式~(\ref{eq:classical_gaussian}) 存在严格的理论适用边界与工程劣势：
\begin{itemize}[leftmargin=2em]
  \item \textbf{理论适用范围受限：} 该公式仅在 $\varepsilon \in (0, 1]$ 的窄区间内被严格证明满足定义；当 $\varepsilon > 1$ 时，该闭式界限不再构成严格的充分条件。
  \item \textbf{噪声方差过度放大（Over-conservatism）：} 经典充分条件公式基于高斯尾部放缩，计算得到的标准差 $\sigma$ 存在松弛。在样本量敏感的企业分析场景下，这种不必要的噪声膨胀将直接损害下游业务的决策可用性。
\end{itemize}

\subsection{解析高斯机制校准 (Analytic Gaussian Calibration)}
为了在保证数学证明严格性的同时最大化统计效用，本项目在底层机制中全面采用 Balle 与 Wang (ICML 2018) 提出的\textbf{解析高斯机制（Analytic Gaussian Mechanism）}~\cite{balle2018improving}。

\begin{techbox}{解析高斯校准定理 (Balle \& Wang, 2018)}
对于给定全局 $L_2$ 敏感度 $\Delta_2$ 以及目标差分隐私参数 $(\varepsilon, \delta)$，高斯扰动机制满足 $(\varepsilon, \delta)$-DP 当且仅当隐私损失随机变量的尾部累积概率函数满足：
\begin{equation}
  \label{eq:analytic_tail}
  \delta(\sigma; \varepsilon, \Delta_2) = \Phi\left(\frac{\Delta_2}{2\sigma} - \frac{\varepsilon\sigma}{\Delta_2}\right) - e^\varepsilon \Phi\left(-\frac{\Delta_2}{2\sigma} - \frac{\varepsilon\sigma}{\Delta_2}\right) \le \delta
\end{equation}
其中 $\Phi(x) = \frac{1}{\sqrt{2\pi}} \int_{-\infty}^x e^{-t^2/2} dt$ 为标准一维正态分布的累积分布函数（CDF）。
\end{techbox}

在工程实现中，系统将式~(\ref{eq:analytic_tail}) 的等号条件构造为关于噪声标准差 $\sigma$ 的一维严格单调函数：
\begin{equation}
  g(\sigma) = \Phi\left(\frac{\Delta_2}{2\sigma} - \frac{\varepsilon\sigma}{\Delta_2}\right) - e^\varepsilon \Phi\left(-\frac{\Delta_2}{2\sigma} - \frac{\varepsilon\sigma}{\Delta_2}\right) - \delta = 0
\end{equation}
采用 Brent 单变量求根算法（\texttt{scipy.optimize.}\allowbreak\texttt{brentq}），下界为 $10^{-5}\Delta_2$，上界取 $\max(2.5b,100\Delta_2)$，其中 $b=\Delta_2\sqrt{2\ln(1.25/\delta)}/\min(\varepsilon,1)$；绝对及相对容差均为 $10^{-12}$。解析条件适用于任意 $\varepsilon>0$，但代码采用浮点求根，不宣称极端参数下的数值完备性。最小方差结论限于给定敏感度的各向同性高斯机制族，不是所有 DP 机制的全局最优。

\subsection{实际参数校准与对比分析}
在基准安全配置（$\varepsilon = 1.0, \, \delta = 10^{-6}$）下，各统计指标的理论敏感度以及分别使用经典公式与解析高斯算法求得的噪声标度 $\sigma$ 对比如表~\ref{tab:sigma_comparison} 所示。

\begin{table}[htbp]
\centering
\small
\caption{基准参数配置下各统计指标的高斯噪声标度校准对比表 ($\varepsilon=1.0, \delta=10^{-6}$)}
\label{tab:sigma_comparison}
\begin{tabularx}{\textwidth}{Xcccc}
\toprule
\textbf{统计指标} & \textbf{配置} & $\boldsymbol{\Delta_2}$ & $\boldsymbol{\sigma_{\rm classical}}$ & $\boldsymbol{\sigma_{\rm analytic}}$ \\
\midrule
User Prevalence & $k=9$ & $3.0000$ & $15.8964$ & $12.6740$ \\
Normalized Share & $\|x_u\|_2\le1$ & $1.0000$ & $5.2988$ & $4.2247$ \\
Interaction Volume & $C=20$ & $20.0000$ & $105.9761$ & $84.4936$ \\
Success Rate & $C_{\text{rate}}=20$ & $28.2843$ & $149.8728$ & $119.4920$ \\
Top-K Tasks & $k_t=5$（保守） & $2.2361$ & $11.8485$ & $9.4467$ \\
\bottomrule
\end{tabularx}
\end{table}

在本文的基准安全配置（$\varepsilon = 1.0, \, \delta = 10^{-6}$）下，实测表明解析 Gaussian calibration 相对当前采用的经典充分条件公式将各指标的噪声标准差 $\sigma$ 降低约 $20.2711\%$（噪声方差降低约 $36.4330\%$），显著增强了小样本聚合单元的抗畸变能力。在其他 $(\varepsilon, \delta)$ 参数设置下，降噪比例会随目标参数组合动态变动，系统不对其他参数组合做统一的百分比承诺。

表中数值由 \texttt{revision/recalculate.py} 调用项目函数并以独立 CDF 求根交叉核对。单位敏感度经典系数为 $5.2988025269$，解析系数为 $4.2246788893$。降幅分别按 $1-\sigma_a/\sigma_c$ 与 $1-(\sigma_a/\sigma_c)^2$ 计算。当前 $T=3$ 的 Top-K 紧界为 $\sqrt3$，但实现保守使用 $\sqrt5$，故表格保留实际校准尺度。
